\head{第十二版新核查：第0章21.5--21.9}
\source{EGA IV第一部分，分册印刷第164--182页的相关部分／上传PDF第163--181页；正式勘误见第四卷第四部分PDF第344页}
本轮从可恢复的第十一版工程接续。展开该工程所得单文件源码与用户已收到的第十一版源码逐字一致，
其连续终点为（21.4.8）。本轮新增21.5--21.9的全部正文，止于（21.9.8）及其证明；
第0章§21至此连续译完。§22章首仅用于确认接续位置，未计入本轮译文。

新增正文保留原著的组织、编号、证明路线与记号。在能够给出明确论证的地方，
校正后的表述写入正文，原文差异放在带L12标识的脚注。
原作者的正式勘误、译者发现的原断言反例、证明步骤补足及边界提醒分别列出。
\textbf{下面各项性质不同，不能把脚注数或表格行数当作错误定理的数量。}

\section*{范围与分类}
\addcontentsline{toc}{section}{范围与分类}
\begin{longtable}{@{}p{.17\linewidth}p{.20\linewidth}p{.54\linewidth}@{}}
\toprule 标识&位置&性质及处理\\\midrule\endhead
L12-E01&21.5.2，c)&原条件有反例；补入一次分次映射非零，给出完整修订论证。\\
L12-E02&21.5.3，ii&原导子扰动未必乘法相容；改用平方零商。唯一性判据不变。\\
L12-A01&21.7.1&落实作者的正式引文勘误。\\
L12-T01--T08&相应脚注&分别核对内部引文、标量与基域记号、两处明确的符号误植，以及作者的语法性勘误。\\
L12-G01--G14中的实际列项&相应脚注&展开原文使用的提升、滤过、有限性、自然性与反例计算；没有补造原著未写的章节。\\
L12-N01--N08&相应脚注&区分边界条件、有限秩前提、完备与普通微分模、分式域与迭代Laurent域、系数域选择。\\
\bottomrule
\end{longtable}
脚注标识为稳定的追踪标识，不要求每个字母序列连续。
旧版H8-P01等问题没有因本轮译校而自动解决；全部历史报告原样保留并另标历史范围。

\audit{L12-E01}{MacLane判据的条件c)缺少一次非退化性}{原文按字面假设不成立；补条件后原等价框架可保留}
\oldauditsrc{$0_{\mathrm{IV}}$（21.5.2），印刷第164--166页，尤其第165页的条件c)}
设 $A,B$ 为完备离散赋值环，极大理想为 $\mathfrak m,\mathfrak n$，
$u:A\to B$ 满足 $Bu(\mathfrak m)=\mathfrak n$。
记 $K=A/\mathfrak m$、$L=B/\mathfrak n$，以及
$u_1:A/\mathfrak m^2\to B/\mathfrak n^2$。
原条件c)只要求另一同态 $u'_1$ 为局部同态且
$\gr_0(u'_1)=\gr_0(u_1)$，随后要求存在自同构 $w_1$，使
$\gr_0(w_1)=\id_L$、$u'_1=w_1u_1$。

\lead{反例。}
取任意域 $k$，置
\[
 A=B=k[[t]],\qquad u=\id,
 \qquad A/\mathfrak m^2=B/\mathfrak n^2=k[t]/(t^2).
\]
在该商环上定义
\[
 u'_1(a+bt)=a\qquad(a,b\in k).
\]
这确实是保单位的环同态；其逆像极大理想为 $(t)$，故是局部同态，
在剩余域 $k$ 上是恒等映射。可是 $u'_1(t)=0$，所以不是单射，更不是自同构。
若有 $u'_1=w_1u_1$，因 $u_1=\id$，就必须 $w_1=u'_1$，矛盾。
与此同时，原条件a)“$L/K$可分”在这里显然成立，因为 $L=K=k$。
因此原文的 $\mathrm a)\Rightarrow\mathrm c)$ 确实失效。

\lead{校订后的条件。}
在c)中补入
\[
 \gr_1(u'_1)\ne0.
\]
这是必要的：$u_1$ 把均匀化参数的类映成非零类，
而自同构在 $\mathfrak n/\mathfrak n^2$ 上必须诱导同构。
它也恰好补足原证明的需要。若 $L/K$ 可分，$B$ 是Cohen $A$-代数，
Cohen提升给出 $h:B\to B/\mathfrak n^2$，使 $h\circ u=u'_1\circ q$，
其中 $q:A\to A/\mathfrak m^2$ 为商映射，且 $h$ 在剩余域上为恒等。
$h$ 是局部同态，故 $h(\mathfrak n^2)=0$，从而诱导 $w_1:B/\mathfrak n^2\to B/\mathfrak n^2$。

因为 $\gr_0(w_1)=\id_L$，$\gr_1(w_1)$ 是一维 $L$-向量空间
$\mathfrak n/\mathfrak n^2$ 的线性自映射。
选择 $A$ 的均匀化参数 $\pi$，由 $u'_1(\bar\pi)\ne0$ 得
$\gr_1(w_1)\ne0$，因此它是同构。
$w_1$ 在有限滤过的零次与一次均为同构，于是自身为同构。
这证明a)推出校订后的c)。

原c')要求 $\gr_1(u'_1)=\gr_1(u_1)$，故自动满足新增的非零条件。
原证明中c')推出a)的导子延拓论证不受影响：对任意导子 $D:K\to L$，
将其看成 $A/\mathfrak m^2\to\mathfrak n/\mathfrak n^2$ 的导子，形成 $u'_1=u_1+D$。
满足c')的自同构差 $w_1-\id$ 在 $\mathfrak n/\mathfrak n^2$ 上为零，
因而下降为 $L\to L$ 的导子，且延拓 $D$。
（20.6.3）遂给出 $L/K$ 可分。
因此正文保留
\[
 \mathrm c)\Longleftrightarrow\mathrm c')\Longleftrightarrow\mathrm a)
 \Longrightarrow\mathrm b)\Longrightarrow\mathrm b').
\]
若 $A$ 为Cohen环，b')推出c')仍由（19.8.6，i）提升给定 $u'_1$ 得到。
在混合特征的非域Cohen环情形，均匀化参数可取整数 $p$；保单位同态必须保持 $p$，
所以此时非零的一次映射往往已经由具体结构保证。
但原命题先处理一般完备离散赋值环，不能因此省去一般条件c)的限定。

\lead{本项的界限。}
这不是对MacLane可分性理论的否定，也不是作者已发表勘误的转述。
它是对这一印刷表述的具体反例和最小非退化条件修订。
正文同时将证明中误引的（20.6.5）改为实际使用的（20.6.3），记为L12-T01。

\audit{L12-E02}{Cohen提升唯一性证明中的平方零商}{原中间步骤缺少必要条件；正文用另一商补证，结论不变}
\oldauditsrc{$0_{\mathrm{IV}}$（21.5.3，ii），印刷第166页}
原文设 $W$ 为Cohen环，$C$ 为完备诺特局部环，$0\ne\mathfrak J$ 包含于 $C$ 的极大理想；
在固定 $W\to C/\mathfrak J$ 的局部同态后讨论提升的唯一性。
当 $\Omega^1_K\ne0$（$K=W/\mathfrak m$）时，原证明取
$D:K\to\mathfrak J/\mathfrak m\mathfrak J$ 非零，
并直接在 $C/\mathfrak m\mathfrak J$ 中用 $u_1+D'$ 扰动提升。

对于环同态 $u$ 与导子 $D'$，乘法差为
\[
 (u(x)+D'(x))(u(y)+D'(y))-(u(xy)+D'(xy))=D'(x)D'(y).
\]
因此，导子加法自动给出环同态，需要导子的值域理想平方为零。
原假设没有给出 $\mathfrak J^2\subseteq\mathfrak m\mathfrak J$。

\lead{原步骤失败的具体例子。}
取特征零的Cohen域 $W=\mathbf Q(t)$，$C=W[[x]]$，$\mathfrak J=(x)$，
$u:W\hookrightarrow C$ 为常数嵌入。此时 $\mathfrak m=0$。
导子 $D'(a)=x\,\partial a/\partial t$ 的值属于 $\mathfrak J$。
令 $v_1=u+D'$，则
\[
 v_1(t^2)=t^2+2tx,
 \qquad v_1(t)^2=(t+x)^2=t^2+2tx+x^2.
\]
因为 $x^2\ne0$，$v_1$ 不是环同态。
这个例子针对原证明的扰动步骤，不反驳所要证明的唯一性判据。

\lead{修补。}
固定一个已有局部提升 $u:W\to C$，以 $\mathfrak r$ 记 $C$ 的极大理想。
设
\[
 I=\mathfrak m\mathfrak J+\mathfrak J^2,
 \qquad Q=\mathfrak J/I.
\]
这里 $\mathfrak m\mathfrak J$ 使用 $u$ 所定义的 $W$-作用。
则 $Q$ 是 $C/I$ 中的平方零理想，且被 $\mathfrak m$ 零化，因而为 $K$-向量空间。
$u$ 局部且 $\mathfrak J\subseteq\mathfrak r$，所以
$I\subseteq\mathfrak r\mathfrak J$。
由于 $C$ 诺特，$\mathfrak J$ 为有限型 $C$-模；若 $Q=0$，则
$\mathfrak J=I\subseteq\mathfrak r\mathfrak J$，由Nakayama引理得 $\mathfrak J=0$，矛盾。
故 $Q\ne0$。

由 $\Omega^1_K\ne0$，存在非零 $K$-线性映射 $\Omega^1_K\to Q$，
即非零导子 $D:K\to Q$。与 $W\to K$ 复合得到 $D':W\to Q$。
在 $C/I$ 中令 $v_1=u_1+D'$。因为 $Q^2=0$，上述乘法差为零，
所以 $v_1$ 确为环同态。它与 $u_1$ 在 $C/\mathfrak J$ 中相同、在剩余域上相同，
故仍是局部同态；又 $D'\ne0$，所以 $v_1\ne u_1$。
$C$ 为完备诺特局部环，$I$ 包含于其极大理想，故（19.8.6，i）将 $v_1$ 提升到 $v:W\to C$。
$v$ 是原给定映射的另一提升，且 $v\ne u$。这完成了必要性方向的修补。

反方向也须区分普通微分模与完备微分模。若 $\Omega^1_K=0$，在混合特征情形
$\mathfrak m=pW$，关系 $d(p)=0$ 与微分正合列给出
$\Omega^1_W=\mathfrak m\Omega^1_W$。
由此得到的是所有 $\mathfrak m$-进制商均为零，故
$\widehat\Omega^1_W=0$；不能对未证明有限型的 $\Omega^1_W$ 直接使用Nakayama引理。
形式非分歧遂给出各级提升的唯一性，完备分离性给出 $C$ 中的唯一性。

\audit{L12-G04}{“上同调相同”与“诱导同态为同构”}{补足原证明使用的自然性，不改变缺陷加法性}
\oldauditsrc{（21.6.4），印刷第168--169页}
对短正合复形 $0\to T_1\to T_2\to T_3\to0$，
仅有 $H^i(T_1)\simeq H^i(T_2)$ 的抽象同构，并不能推出 $T_3$ 无上同调。
所需条件是给定映射 $H^i(T_1)\to H^i(T_2)$ 本身为同构。
一个简单的警告是对集中在零次的一维向量空间取零映射；两端对象相同，映射却不是同构。

在原证明（i）中，两行复形的唯一可能非零上同调都通过（21.6.1.1）典范识别为
\[
 \Upsilon_{k''/k'/k}\otimes_{k''}L.
\]
由域嵌入和非完美性模的自然性，原图给出的上同调映射正是这一空间的恒等映射。
故可用长正合列得第三行正合。
在（ii）中，共同的上同调空间为 $\Upsilon_{M/L/K}$，诱导映射同样是恒等。
正文把这两个必要核对写出，保留缺陷的两条加法公式。

\audit{L12-G06}{Cartier等式中的有限性先于维数相减}{保留定理；补出有限生成归纳的逻辑次序}
\oldauditsrc{（21.7.1）--（21.7.3），印刷第170--171页}
Cartier等式为
\[
 \dim_L\Omega^1_{L/K}-\dim_L\Upsilon_{L/K}=\trdeg_K L
\]
其中 $L/K$ 为有限生成域扩张。
原证明利用六项正合列以及单生成扩张归纳。
必须同时知道参与相减的空间有限维，不能先写未定义的 $\infty-\infty$。

对 $L=K(x)$，若 $x$ 超越，则可分性给出 $\Upsilon=0$，微分模秩为1；
若有限可分，则两者均为零；若代数不可分，先经过最大可分中间域，
再拆成有限个次数为 $p$ 的纯不可分扩张，在高度一情形由（21.4.8）得两空间同秩且有限。
对有限多个生成元归纳时，六项正合列中其余空间已经有限维，
从而同时推出新一步的有限维性与Euler特征等式。
正文遵循这一顺序。作者已将此处引用（20.6.15.1）正式改为（20.6.17.1）；
这项作者勘误标为L12-A01，不作为译者新发现。

\audit{L12-N03}{逆极限与扩张标量的使用边界}{只证明所需单射性，不宣称任意交换逆极限和张量积}
\oldauditsrc{（21.8.3）的 $\mathrm c)\Rightarrow\mathrm b)$，印刷第173页}
设 $W\to\varprojlim W_\alpha$ 单射，所有空间均在域 $K$ 上。
原证明需要，对任意 $K$-向量空间 $V$，映射
\[
 W\otimes_KV\longrightarrow\varprojlim(W_\alpha\otimes_KV)
\]
仍单射。这里并不需要或声称一般同构
$(\varprojlim W_\alpha)\otimes_KV\simeq\varprojlim(W_\alpha\otimes_KV)$。

若 $V$ 有限维，选取有限基，所需结论只是有限直和保持单射。
对一般 $V$，任一张量只涉及有限多个 $V$ 的向量，故位于某个有限维子空间 $V_0$ 的张量积中。
向量空间张量积是正合的，所以每个
$W_\alpha\otimes V_0\to W_\alpha\otimes V$ 都单射；
逆极限保左正合，故逆极限后的映射也单射。
由此即可检测每个非零张量，不必假设任何无限交换律。

\audit{L12-T06}{无限秩反例中的基域应为 $K^p$}{原文基域误植；原作者反例经核对后保留}
\oldauditsrc{（21.8.7，ii），印刷第175页}
原文说构造一列“$K$的有限扩张 $M_n$”，随后却定义
$M_0=K^p$、$M_n=K^p(a_1x,\ldots,a_nx)\subseteq K$。
两者不相容；正确的是 $K^p$ 在 $K$ 中的一列有限扩张。

为核验原构造，取 $k_0=\mathbf F_p$、$[K:K^p]$ 可数无限，
取 $K=K_0\supsetneq K_1\supsetneq\cdots$ 且 $\bigcap K_n=K^p$，
并保证每个 $K_n/K^p$ 仍为无限次扩张。
固定 $x\in K\setminus K^p$。归纳选
\[
 a_{n+1}\in K_{n+1}\setminus M_n(x),\qquad
 M_{n+1}=M_n(a_{n+1}x).
\]
若 $F=M_n$，则 $x\notin F$、$a_{n+1}\notin F(x)$，且这两个元素的$p$次幂都属于 $K^p\subseteq F$。
所以 $[F(x,a_{n+1}):F]=p^2$。
但 $[F(a_{n+1}x):F]\le p$。
若 $x\in F(a_{n+1}x)$，就也有 $a_{n+1}=(a_{n+1}x)/x\in F(a_{n+1}x)$，矛盾。
于是始终 $x\notin M_n$。

置 $M=\bigcup M_n$、$L=M^{1/p}$。因 $K^p\subseteq M$，$L$ 包含 $K$，且 $L^p=M$。
对每个 $n\ge1$，$x=(a_nx)/a_n\in K_n(M)=K_n(L^p)$，
但 $x\notin M=L^p$。因此
\[
 \bigcap_nK_n(L^p)\supsetneq k_0(L^p),
 \qquad \bigcap_nK_n(K^p)=k_0(K^p).
\]
这验证了原作者要说明的“不能删去有限秩假设”。本版并不把该反例宣称为新发现。

\audit{L12-N04}{多变量形式幂级数的分式域及分母统一}{记号和证明步骤澄清}
\oldauditsrc{（21.8.8）及子引理（21.8.8.2）、（21.8.8.3），印刷第175--176页}
本处 $k((T_1,\ldots,T_r))$ 表示
$\operatorname{Frac}(k[[T_1,\ldots,T_r]])$，不是迭代Laurent级数域。
两个域在多变量情形不可随意替换。

子引理（21.8.8.3）的证明中，固定 $f$ 属于所有 $k_\alpha((T))$，
先找 $g\in k_\gamma[[T]]\setminus\{0\}$，使 $gf\in k_\gamma[[T]]$。
若 $k^q\subseteq k_0=\bigcap k_\alpha$，将 $g$ 换为 $g^q$ 后，
分母属于 $k_0[[T]]$，同时 $g^qf=g^{q-1}(gf)$ 仍是 $k_\gamma$-系数的幂级数。
对共尾的 $k_\alpha\subseteq k_\gamma$，应用前一子引理于 $g^qf$，
得到它属于每个 $k_\alpha[[T]]$，因而属于 $k_0[[T]]$。
所以 $f\in k_0((T))$。这一步针对乘过统一分母后的元素，而不是无条件断言 $f$ 本身是幂级数。

\audit{L12-G11}{完备微分模中密度、连续性与有限性的区分}{补明适用条件，保留命题}
\oldauditsrc{（21.9.2）、（21.9.4）、（21.9.5），印刷第176--178页}
（21.9.4）使用“导子在稠密多项式子环上为零，故在形式幂级数环上为零”。
连续映射在稠密子集上为零，要推出处处为零，目标应是分离的。
正文限定在分离完备测试模中使用这一论证；这些测试模已经足以使用完备微分模的泛性质。
此外，$A$ 在 $B_1=k[[T_1^p,\ldots,T_r^p]]$ 上有限型，故 $\Omega^1_{A/B_1}$ 是有限型 $A$-模；
由于 $A$ 为完备诺特局部环，这个模分离完备，因而不用再取额外完备化。

（21.9.5）中，延拓到有限可分分式域扩张的导子需要具有共同的有限型值域并连续。
选取有限个 $A_0$-代数生成元 $a_j$，令
\[
 P=A+\sum_{i,j}A\,D_i(a_j)\subseteq E.
\]
Leibniz公式表明每个 $D_i(A)\subseteq P$；$P$ 为有限型 $A$-模，也为有限型 $A_0$-模。
若 $\mathfrak m_0=(T_1,\ldots,T_r)$，则
\[
 D_i(\mathfrak m_0^nA)\subseteq\mathfrak m_0^{n-1}P\qquad(n\ge1),
\]
所以这些导子连续。其在参数上的值 $D_i(T_j)=\delta_{ij}$ 给出所需线性无关性。
正文把原文这一有限值域和连续性的核对写出。

还需注意，完备微分模的正合性并非来自“完备化总是正合”。
在（21.9.2）中，连续左逆保证像为闭子空间，再结合原有的稠密性才得到像等于核；
在（21.9.5）中，则使用诺特有限型模子模闭的性质。两条理由不能混为一般交换律。

\audit{L12-T08}{原作者的无限 $p$-秩例子：原印误与文本识别错误分开}{校正 $p$-基符号；保留并补足原反例的验证}
\oldauditsrc{（21.9.7），印刷第181页／PDF第180页}
原页写 $k_0=\mathbf F_p(X_n)_{n\ge0}$，接着把其在 $\mathbf F_p$ 上的$p$-基写成
$(X_n^p)$。这不成立：每个 $X_n^p$ 都在 $k_0^p$ 中，
其微分为零，不能成为$p$-基。正确的是 $(X_n)$。
本轮放大原页确认该$p$次幂确为原印刷符号，故正文校正并在脚注记录。

\textbf{另一方面，原页在“否则”后正确印着 $y^p$，不是 $y$。}
文本提取曾丢掉上标，不能把这种识别差异算作原作者错误。
正文保留正确的 $y^p$，并另补证明其不属于相关复合域。

具体设
\[
 K_0=k_0=\mathbf F_p(X_0,X_1,\ldots),\quad
 L_0=k_0((T)),\quad M_0=k_0((T^p)),
\]
\[
 f=\sum_{n\ge0}X_nT^{np},\qquad E=L_0(y),\quad y^p=f.
\]
任一 $k_0(L_0^p)$ 的元素只用到有限多个 $X_i$ 作为额外系数，
故属于某个
$k_0^p(X_i:i\in I)((T^p))$，其中 $I$ 有限。
$f$ 的所有系数 $X_n$ 不可能同时属于这样一个有限生成的系数子域，
由形式级数系数的唯一性，
\[
 f\notin k_0(L_0^p).
\]
因此 $d_{L_0/k_0}f\ne0$。

由形式幂级数分式域的可分性，$L_0/k_0$ 可分。
在绝对微分模的标准表示
\[
 \Omega^1_{E/\mathbf F_p}
 \simeq
 \left((\Omega^1_{L_0/\mathbf F_p}\otimes_{L_0}E)\oplus E\,dy\right)/(E\,df)
\]
中，$df$ 的相对 $k_0$ 微分非零，故直线 $E\,df$ 与
$\Omega^1_{k_0/\mathbf F_p}\otimes_{k_0}E$ 的像交为零。
后者在 $\Omega^1_{L_0/\mathbf F_p}\otimes E$ 中原本单射，
于是其到 $\Omega^1_{E/\mathbf F_p}$ 的映射仍单射。
由（20.6.3），$E/k_0$ 可分，所以 $d(E/K_0,k_0)=0$。

又有 $E^p\subseteq M_0$、$[L_0:M_0]=p$、$[E:L_0]=p$，
因此 $[E:M_0]=p^2$，$T,y$ 为 $E/M_0$ 的$p$-基，
$\dim_E\Omega^1_{E/M_0}=2$。
Cartier等式给出 $\dim_E\Upsilon_{E/M_0}=2$。
由于基域参数取 $M_0$ 本身，$d(E/M_0,M_0)=2$。
原式中的差于是为
\[
 \delta=d(E/M_0,M_0)-1=1\ne0.
\]
这核验了原作者反例所需结论：不能把（21.9.6，ii）推广到任意无限 $p$-秩的 $k_0$。

\audit{L12-N06}{有限性前提和系数域选择}{明确继承条件，不扩大为未经证明的新定理}
\oldauditsrc{（21.9.6）及（21.9.8），印刷第178--182页}
（21.9.6）的维数差是在 $[K_0:k]<\infty$ 及相应微分空间有限秩的前提下写出的。
把 $k$ 换为更小的 $k'$ 时，正文明确仍需满足所用表达式的有限性前提。
这不是宣布原命题有新的反例，而是防止离开讨论范围后把无穷维的差当作整数。

（21.9.8）先指定 $K_0$ 为 $A$ 的剩余域的子域，随后在证明中把 $K_0$ 提升到 $A$。
为了让 $\Omega^1_{(A/\mathfrak p)/k}$ 的结构从一开始明确，正文先固定一个与剩余域嵌入相容的
$k_0$-嵌入 $K_0\to A$。它由可分性、形式光滑性和 $A$ 的完备性存在。
一般没有唯一性，故本版不把这个选择说成典范。
原证明对任一这样的选择均适用，所得不等式及适当子域上取等号的结论保持不变。

\section*{其余排印、引文与核查记录}
\addcontentsline{toc}{section}{其余排印、引文与核查记录}
\begin{longtable}{@{}p{.22\linewidth}p{.68\linewidth}@{}}
\toprule 位置&处理\\\midrule\endhead
21.5.3末句&作者将法文性数配合ramifiée改为ramifié；中文“形式非分歧”不受影响。\\
21.6.3&原引20.6.17，ii无对应分项；改为20.6.21，ii中的自然性，并保持不同标量域间的双同态解释。\\
21.7.5&同调识别指向20.6.5.2，两项复形指向20.6.5，末尾$F_\bullet(C/A)$指向20.6.26；不是重新定义这些复形。\\
21.8.3&高度一扩张的非完美性模基对应21.4.8；非零微分的判据对应21.4.6。\\
21.9.6证明&印刷第179页首个维数等式中的$\Omega^1_{M/M_0}$方向颠倒，改为$\Omega^1_{M_0/M}$。后一模的标量域为$M_0$，且正是21.7.3所给项。\\
21.9.6.7&$k_0/k_0^p$的有限基扩到更大的系数域后只需生成，不声称仍线性无关。\\
\bottomrule
\end{longtable}

\section*{来源、外部交叉检查与未决事项}
\addcontentsline{toc}{section}{来源、外部交叉检查与未决事项}
翻译底本是用户上传的1964年法文原版，不是从二手综述重构。
本轮阅读了印刷第164--182页的连续内容，并对公式、图和小字号位置查看原页图像；
作者正式勘误以1967年第四部分所载《清单三》原页为依据。

可分域扩张与形式光滑性的背景对应关系，另参照Stacks Project的
\texttt{Tag 031U}（\url{https://stacks.math.columbia.edu/tag/031U}，本轮核对日期：2026年11月16日）。
这一交叉检查不是对上述具体原文错误的既有勘误归属证明；
L12-E01、L12-E02等结论均以本报告给出的具体计算或修补为依据。

本轮没有重新审定全部旧译文，没有把已知待核事项静默删除。
H8-P01的原版较弱拓扑范围仍待核；第四卷主体及第五卷预备笔记余文仍未连续译入。
本报告的历史部分保存旧判断，不表示本轮重做了全部反例与证明。
实际构建、页数、逐页比较和视觉抽查结果记录在工程的\texttt{checks/}中，
本报告不预先将技术检查写成已通过，也不把编译成功当作数学正确性的充分证明。
